\section{Special plates}

\subsection{Imported vehicles (1997 and 1998)}\label{sec:importados}

Plates of normal format and colours from \plate{00-00-KA} to \plate{00-00-KF} were
reserved for second-hand imported vehicles registered during 1997.

For vehicles imported by private individuals in 1997 and 1998, the first letter
of the letter group was also a \plate{K}; the highest combination issued is said
to have been \plate{KG}.\footnote{Portuguese Wikipedia says every group from \plate{KA} to
\plate{KZ} was assigned to imports in 1997; the two sources disagree.}

\platefig{import-8949kc}{Imported vehicle: \plate{89-49-KC}.}

\subsection{New vehicles not yet registered}

Red plates with white characters, always 220~mm $\times$ 330~mm, for new vehicles
not yet registered. The top shows a sequential number and the bottom the name (or
abbreviation) of the vehicle's registered operator, manufacturer or dealer.

\platefig[0.3\linewidth]{new-6019-santagri}{New vehicle not yet registered.}

\subsection{Competition vehicles}

Red plates with black characters, similar to those of industrial machines and
trailers, but with the year and month of first registration.

\platefig{competition-39th48}{Competition vehicle: \plate{39-TH-48}.}

\subsection{Vehicles for export}

Yellow plates with black characters and the sequence \plate{00000-X}, where
\plate{X} identifies the customs office where the vehicle was registered. On the
yellow band, the letters \plate{EXP} are added above the date.

\gen{export}

\platefig[0.35\linewidth]{export-00573l}{Vehicle for export: \plate{00573-L}.}

\subsection{Industrial machines and trailers}

Red plates with black characters. On the yellow band, instead of the date, a
letter shows the machine's or trailer's traffic class (classes \plate{A} to
\plate{H}). Their registration is governed by Decree-Law 107/2006 (RAMMI). The
plates use the normal car sequence.

\platepair{industrial-gp1234}{Industrial machine: \plate{GP-12-34}, class \plate{A}.}%
          {industrial-48xj68}{Industrial machine: \plate{48-XJ-68}, class \plate{A}.}

\subsection{Agricultural tractors}

Agricultural tractors and their trailers have plates with a green background
(Regulatory Decree 13/98).

\subsection{National Republican Guard (GNR)}

Plates of normal format and colours with the sequence \plate{GNR X-0} to
\plate{GNR X-0000} are reserved for the National Republican Guard, where
\plate{X} is a letter showing the vehicle's type or function (see
appendix~\ref{sec:gnr}).

\subsection{Military forces}

\begin{description}
  \item[Portuguese Army] Plates of normal format and colours with the sequences
    \plate{MG-00-00}, \plate{ME-00-00} and \plate{MX-00-00}. Most still use the
    black plates with white letters.
  \item[Portuguese Navy] Sequence \plate{AP-00-00}.
  \item[Portuguese Air Force] Sequence \plate{AM-00-00}.
  \item[Allied forces] A layout similar to ordinary Portuguese plates, but with
    the letter \plate{U} on the euroband instead of \plate{P}. Used mainly by US
    military personnel at Lajes Air Base, on Terceira in the Azores, since 2005.
\end{description}

\platepair{army-mx5531}{Army: \plate{MX-55-31}.}%
          {allied-b589an}{Allied forces: \plate{B5-89-AN}, with \plate{U} on the euroband.}

\subsection{Vehicles of the President of the Republic}

Perhaps the most distinctive of the special plates, used on every car assigned to
the President of the Republic: a black background with the national coat of arms
in the centre, flanked by the letters \plate{P} and \plate{R} in white or silver.

\platefig[0.3\linewidth]{president}{Plate of the President of the Republic's vehicles.}

\subsection{Diplomatic corps}\label{sec:diplomaticos}

White plates with red characters and no euroband, with the sequence
\plate{000-XX000}, reserved for diplomats and staff on international missions.
Decree-Law 54/2005 (art.~6) extends the red characters, dashes and border to bodies
covered by the Protocol on the Privileges and Immunities of the European
Communities.
The first group of digits identifies the country or international organisation
(see appendix~\ref{sec:codigos-diplomaticos}), and \plate{XX} can be:
\begin{itemize}
  \item \plate{CD}: Diplomatic Corps;
  \item \plate{CC}: Consular Corps;
  \item \plate{FM}: official on an international mission.
\end{itemize}

\platepair{diplomatic-cd}{Diplomatic Corps: \plate{206-CD444}.}%
          {diplomatic-cc}{Consular Corps: \plate{024-CC410}.}
